\subsection{Positive hermitian forms}

DEFINITION 2. --- Let E be a vector space over the field K. A hermitian form f on E is said to be positive if \(f(x, x) \geqslant 0\) for all \(x \in E\) .

It is clear that hermitian forms on a vector space E form a vector space over the field R (but not over the field C, when K is C): in this space the positive hermitian forms constitute a pointed convex proper cone (II, p.~10) as a result of def. 2 and prop. 1.

PROPOSITION 2. --- If \(f\) is a positive hermitian form, we have

\[
\left| f (x, y) \right| ^ {2} \leqslant f (x, x) f (y, y)\tag{7}
\]

for every x and y in E (Cauchy-Schwarz inequality).

First assume that we have \(f(y, y) \neq 0\) . For every \(\xi \in \mathbf{K}\) , we have

\[
f (y, y) f (x + \xi y, x + \xi y) \geqslant 0
\]

which can be written as

\[
f (x, x) f (y, y) - | f (x, y) | ^ {2} + (\xi f (y, y) + \overline {{f (x , y)}}) (\bar {\xi} f (y, y) + f (x, y)) \geqslant 0.
\]

Replacing \(\xi\) by \(-\overline{f(x,y)} / f(y,y)\) in this inequality, we get (7). If \(f(x,x)\neq 0\) , we argue similarly.

Finally, if \(f(x, x) = f(y, y) = 0\) , we have \(f(x + \xi y, x + \xi y) \geqslant 0\) for all \(\xi \in \mathbf{K}\) , which can be written as

\[
\xi f (x, y) + \overline {{{{\xi f (x , y)}}}} \geqslant 0.
\]

Replacing \(\xi\) by \(-\overline{f(x,y)}\) in this inequality, we get \(-2|f(x,y)|^2 \geqslant 0\) , and therefore \(f(x,y) = 0\) ; we again get (7) in this case.

COROLLARY 1. --- If \(f\) is a positive hermitian form, the set \(\mathbf{N}\) of all \(x \in \mathbf{E}\) such that \(f(x, x) = 0\) coincides with the vector subspace of all \(x \in \mathbf{E}\) such that \(f(x, y) = 0\) for all \(y \in \mathbf{E}\) .

COROLLARY 2. --- For a positive hermitian form to be separating, it is necessary and sufficient that the relation \(x \neq 0\) implies \(f(x, x) > 0\) .

This follows immediately from cor. 1.

For every positive hermitian form f on E, the separating hermitian form associated with \(f(\mathrm{V}, \mathrm{p.~2})\) is evidently a positive hermitian form on E/N.

PROPOSITION 3. --- Let f be a positive hermitian form on E. Put

\[
p (x) = f (x, x) ^ {1 / 2}
\]

for all \(x \in \mathbb{E}\) . Then \(p\) is a semi-norm on \(\mathbb{E}\) , and is a norm if and only if \(f\) is separating. It is enough to prove the inequality \(p(x + y) \leqslant p(x) + p(y)\) . But we have

\[
f (x + y, x + y) = f (x, x) + f (y, y) + f (x, y) + \overline {{f (x , y)}}
\]

and, by Cauchy-Schwarz inequality

\[
\begin{array}{c} f (x + y, x + y) \leqslant f (x, x) + f (y, y) + 2 (f (x, x) f (y, y)) ^ {1 / 2} \\ = \bigl (f (x, x) ^ {1 / 2} + f (y, y) ^ {1 / 2} \bigr) ^ {2}. \end{array}\tag{8}
\]

Remarks. --- 1) Suppose \(f\) is positive and separating, and let \(x, y\) be two vectors \(\neq 0\) . The proof of Cauchy-Schwarz inequality shows that, if the two members of (7) are equal, then there exists a scalar \(\xi\) such that \(f(x + \xi y, x + \xi y) = 0\) , hence \(x + \xi y = 0\) , in other words, x and y are linearly dependent; the converse is immediate. The proof of inequality (8) shows that the equality \(p(x + y) = p(x) + p(y)\) is possible only if x and y are linearly dependent; if \(y = \lambda x\) , the preceding equality can be written as \(|1 + \lambda| = 1 + |\lambda|\) , and implies that \(\lambda\) is real and positive.

\begin{enumerate}
\def\labelenumi{\arabic{enumi})}
\setcounter{enumi}{1}
\tightlist
\item
  Let \(f\) be a positive hermitian form on \(\mathbf{E}\) , and let \(\mathbf{E}\) be assigned the semi-norm \(x \mapsto f(x, x)^{1/2}\) ; if \(\dot{f}\) is the positive, separating hermitian form defined on \(\mathrm{E/N}\) associated with \(f\) , then the normed space obtained by assigning the norm \(x \mapsto \hat{f}(\dot{x}, \dot{x})^{1/2}\) to \(\mathrm{E/N}\) is the normed space associated with \(\mathbf{E}\) (II, p.~5).
\end{enumerate}

DEFINITION 3. --- Let E be a vector space over the field K. A semi-norm p on E is said to be prehilbertian if there exists a positive hermitian form f on E such that \(p(x) = f(x, x)^{1/2}\) for all \(x \in E\) .

Observe that for a semi-norm p on E, there exists at most one positive hermitian form f such that \(p(x) = f(x, x)^{1/2}\) for all \(x \in E\) ; this follows from the polarization formulas (V, p.~2).

